\begingroup
\titlespacing*{\section}{0pt}{0.8em}{0.45em}

\section{Conclusion}
This thesis presented VietFlood, a multi-platform flood reporting prototype for citizen reporting, relief awareness, and admin review. The project began from a practical problem: flood information often reaches communities through scattered messages, photos, and location shares. That information can help, but it is hard to organize when many reports arrive at the same time. VietFlood gives that process a clearer path. A citizen can submit a report from the mobile app, attach evidence, provide location details, and let the backend store the report for later review.

The completed prototype connects four main parts. The NestJS backend provides the API gateway, authentication service, reports service, and live tracking gateway. The Expo React Native mobile app handles citizen and relief workflows. The Next.js web dashboard handles admin login, report filtering, evidence preview, and status updates. The shared \texttt{vietflood\_common} package keeps Redis, Cloudinary, and logging utilities in one place so the backend services do not repeat the same infrastructure code.

The implementation shows that the main report loop works as a software system. The API gateway accepts report submissions and evidence files. Cloudinary stores uploaded media. PostgreSQL stores users, refresh tokens, and report data. Redis stores report cache entries. RabbitMQ moves backend requests from the gateway to the auth and reports services. Socket.IO handles live location updates with coordinate validation. This design follows the microservices pattern of separating independent capabilities into owned services while using message brokers for communication \cite{fowlerLewis2014microservices,dragoni2017microservices,thones2015microservices}. These parts do not solve every disaster-response problem, but they prove that the prototype can move a citizen report from mobile submission to backend storage and admin review.

Role separation is one of the clearer outcomes of the system. Citizens create and manage their own reports. Relief users read report data and use mobile screens focused on response awareness. Admin users review reports and update report status. This separation keeps verification and rejection under admin control, while relief users can still inspect reports and locations for field work. That boundary is useful because report status affects trust in the system.

The project also shows the value of building the system across mobile, web, backend, and shared code rather than treating flood reporting as only a mobile form. The mobile app is close to the citizen and relief user. The web dashboard gives admins a larger view for review. The backend holds the rules, storage, messaging, and role checks. The shared library keeps repeated infrastructure code under control. Together, these parts make VietFlood a base for later testing and improvement.

At the same time, VietFlood should be understood as a prototype. It has not yet been tested in a real flood situation, with many users, unstable network connections, and urgent field decisions. The current system demonstrates the technical path, not full operational readiness. Before practical use, it needs field testing, stronger upload rules, better monitoring, stricter live-tracking authentication, and more careful review of token behavior.

\section{Future Works}

Future work should begin with field testing. The system needs feedback from citizens, relief workers, and administrators using real phones and realistic network conditions. That testing would show whether the report form is clear, whether evidence upload is reliable enough, and whether the admin dashboard provides enough context for status review.

The mobile app should also support poor-network conditions better. A future version should let citizens save reports locally, retry submission when the connection returns, and compress evidence files before upload. The app should also make the report state clear: sent, saved, uploading, or failed.

Report review can be improved through duplicate detection, spam handling, and evidence quality checks. Reports with similar coordinates, descriptions, and creation times could be grouped for admin review. Relief workflows could also add assignment, response notes, and route history while keeping verification and rejection under admin control.

Security and operations should be stronger before public deployment. Refresh-token handling, file size limits, file type checks, live-tracking authentication, audit logs, health checks, service alerts, and deployment rollback steps all need more work. The operational work should follow DevOps and deployment automation practices \cite{kim2016devopsHandbook,humble2010continuousDelivery}. If report storage or evidence upload fails, administrators should know quickly through monitoring and alerting. The containerized deployment using Docker Compose is already a base, but observability, alerts, and recovery steps are needed before production use \cite{dockerComposeDocs}.

In later research, VietFlood could also add flood insight analysis. The existing reports already contain category, urgency, severity, coordinates, address fields, status, and timestamps. With enough verified reports, the system could generate simple heat maps, time-based charts, or province-level summaries. These additions should come after the reporting path is stable. Bad input will produce bad insight, so data quality must come first.

\endgroup
